\begin{abstract}
Practitioners routinely craft increasingly verbose prompts under the assumption
that more context yields better responses.
We challenge this assumption with a controlled \emph{saturation experiment}:
for six diverse tasks and seven large language models, we measure response
quality at seven additive prompt-length levels and fit quality-versus-token
curves to identify \emph{saturation points}, the token counts at which
quality reaches 95\% of its asymptote.
Across 5,875 LLM-judge-evaluated responses, we find a consistent
\textbf{structured--open dichotomy}: tasks requiring structured output
(classification, product extraction) exhibit clear saturation curves
(F-test $p < 0.05$ in 7 of the $2 \times 7 = 14$ structured model--task
pairs; mean R$^2 = 0.85$ for significant fits), with classification
saturating at just 38--64 tokens.
Product extraction saturates over a wider range (92--536 tokens), with
stronger models requiring far fewer tokens, illustrating that saturation
point is jointly determined by task complexity and model capability.
In contrast, open-ended or reasoning-heavy tasks (question answering, math
reasoning) show no significant curve improvement over a constant-quality baseline
(0 of 14 open-ended pairs), indicating that quality is bottlenecked by task difficulty
rather than prompt elaboration.
Model capability modulates saturation point within structured tasks: stronger
models require fewer tokens to reach asymptote.
These findings provide actionable guidance: structured-output prompts can be
aggressively trimmed after the saturation threshold, while extra context in
open-ended prompts yields negligible return.
\end{abstract}
